\chapter{Introducción}\label{chap:introduccion}

El aprendizaje automático se utiliza hoy en muchos ámbitos de la ingeniería
informática, como la visión por computador, el procesamiento del lenguaje natural, la
medicina, la industria o la ciberseguridad. Buena parte de estos avances se apoya en el
aprendizaje profundo \parencite{goodfellow2016deep}, cuyo rendimiento depende de la
cantidad y la diversidad de los datos con que se entrena el modelo. Esos datos, además,
no siempre están donde se necesitan: rara vez residen en un único repositorio y, con
frecuencia, no pueden trasladarse sin restricciones a una infraestructura central.

El enfoque habitual ha sido el entrenamiento centralizado, que reúne en un servidor
común los datos de las distintas fuentes. En muchos casos reales, en cambio, esos datos
están repartidos entre dispositivos, usuarios u organizaciones, y centralizarlos es
difícil o poco conveniente. La privacidad lo limita, transferir grandes volúmenes de
información tiene un coste elevado y concentrar datos sensibles en un solo punto añade
riesgo \parencite{li2020federated,kairouz2021advances}. La colaboración entre entidades
sigue siendo deseable, pero compartir los datos en bruto no siempre es una opción.

El aprendizaje federado (\textit{Federated Learning}, FL) responde a esa necesidad:
entrena modelos de forma colaborativa sin sacar los datos de los dispositivos o
entidades donde se generan \parencite{kairouz2021advances}. El modelo se entrena
localmente en cada cliente y, en lugar de los datos, solo viajan las actualizaciones
resultantes, que un servidor central combina para formar un modelo global. El algoritmo
de referencia es \textit{Federated Averaging} (FedAvg), propuesto por
\textcite{mcmahan2017communication}, que organiza el proceso en rondas sucesivas de
distribución del modelo, entrenamiento local y agregación.

Implementar un sistema federado, sin embargo, exige bastante más que entrenar una red
neuronal. Hay que gestionar la comunicación entre servidor y clientes, definir
estrategias de agregación, controlar las rondas, registrar métricas y dejar los
experimentos en condiciones de reproducirse \parencite{bonawitz2019towards}. Para no
rehacer toda esa infraestructura cada vez, han aparecido \textit{frameworks}
especializados que ofrecen componentes reutilizables con los que construir y evaluar
este tipo de sistemas \parencite{bonawitz2019towards,caldas2018leaf}. El problema es que
hay varios y no siempre está claro cuál conviene para un caso concreto, porque cada uno
toma decisiones de diseño distintas y puede comportarse de otra manera ante la misma
carga de trabajo.

Ahí se sitúa este Trabajo de Fin de Grado. Se propone comparar varios \textit{frameworks}
de aprendizaje federado de forma experimental: implementar en todos ellos un mismo
escenario, en condiciones equivalentes, y medir tanto el aprendizaje del modelo
(precisión y pérdida) como el coste en recursos (tiempo de ejecución, memoria, uso de
GPU, potencia y energía estimada), la reproducibilidad y la facilidad de uso. No se
busca decidir qué herramienta es mejor en términos absolutos, sino ofrecer una
comparación razonada y reproducible que aclare sus diferencias prácticas y sirva de
orientación en proyectos parecidos.

\section{Contextualización del problema}\label{sec:contextualizacion}

Durante mucho tiempo, entrenar modelos de aprendizaje automático ha dependido de tener
los datos centralizados. Cuando todos pertenecen a una misma organización y pueden
almacenarse y procesarse sin restricciones, este enfoque funciona bien y simplifica el
control del preprocesamiento y la evaluación del modelo. En muchos sistemas actuales, en
cambio, la situación es distinta. Los datos están repartidos entre varias fuentes y
llevarlos a un mismo lugar suele ser problemático \parencite{kairouz2021advances}.

El motivo no es solo la privacidad. La protección de datos personales o sensibles es una
de las razones de peso para no centralizar, pero hay además limitaciones técnicas y
organizativas. El volumen de datos que generan dispositivos y sensores puede hacer poco
eficiente su transferencia continua; el almacenamiento centralizado encarece la
infraestructura y convierte al servidor en un punto crítico de seguridad; y, en el plano
organizativo, dos entidades independientes pueden no estar dispuestas a compartir sus
bases de datos completas \parencite{li2020federated}.

Frente a este panorama, el aprendizaje federado permite entrenar un modelo de forma
colaborativa sin reunir todos los datos en un mismo sitio. Lo que se comparte son las
actualizaciones del modelo, no los datos originales, y eso lo hace apropiado para
escenarios distribuidos o con datos sensibles. Conviene matizar, eso sí, que no elimina
del todo los riesgos de privacidad \parencite{kairouz2021advances}. Estos fundamentos se
desarrollan en el Capítulo~\ref{chap:fundamentos}.

\section{Estructura de la memoria}\label{sec:estructura}

El resto de la memoria se organiza de la siguiente manera. El
Capítulo~\ref{chap:objetivos} concreta el objetivo general y los objetivos específicos
que guían el trabajo. El Capítulo~\ref{chap:fundamentos} reúne los fundamentos teóricos
del aprendizaje automático y del aprendizaje federado, incluido el algoritmo
\textit{FedAvg}, los tipos de aprendizaje federado y los retos que plantea la
heterogeneidad de datos y clientes. Después se revisa el estado del arte y se justifica
la selección de los \textit{frameworks} comparados. El siguiente bloque describe la
metodología experimental: el diseño del experimento, el control de variables, la
reproducibilidad y las métricas empleadas, junto con el alcance y las limitaciones del
trabajo. Por último se presentan los resultados, su análisis comparativo y las
conclusiones.
